
Trustworthiness evaluations of large language models (LLMs) are used to guide model selection for safety-sensitive applications, with the implicit assumption that in-distribution benchmark rankings predict out-of-distribution behavior. This study tests that assumption directly. Using partial Spearman $\rho$ (MMLU-controlled) across 16 LLMs from the TrustLLM evaluation suite \citep{Sun2024TrustLLM}, we compute cross-split rank stability for fairness (BBQ-Disambig $\rightarrow$ BBQ-Ambig) and adversarial robustness (ANLI R1 $\rightarrow$ R3; OOD robustness). Contrary to the hypothesis that adversarial benchmark construction disrupts rank stability, both fairness ($\rho$ = 0.962, p < 0.0001, N = 16) and adversarial robustness ($\rho_{ANLI}$ = 0.684, p = 0.007; $\rho_{OOD}$ = 0.868, p = 0.0001; N = 13) exhibit substantially positive partial $\rho$ after capability control, with zero rank reversals across all adversarial pairs. Fairness shows marginally stronger stability than adversarial robustness ($\Delta\rho$ = 0.192, Fisher z = 2.265, p = 0.024, N = 13); however, the preregistered effect-size criterion ($\Delta\rho \geq$ 0.200) was not met, and the proposed adversarial disruption mechanism is falsified. Trustworthiness rankings appear to be stable latent model properties, not test-specific artifacts, supporting the predictive use of in-distribution trustworthiness evaluations for deployment decisions.

